\section{Automatic application and the next frontier}\label{sec:automatic}
\subsection{A label-free decision procedure}
Every exclusion receives its input as the complete integer multiplication
tensor. The application order is arithmetic screening, the cyclic-four
central-support criterion when recognized, quadratic parity when its
hypotheses hold, and otherwise the positive-indicator sphericalization
framework. Pseudo-unitary branches use the low-induction bounds followed
by the same finite-indicator systems. Spherical branches exhaust the
dimension-character components. Signed branches exhaust the bounded
associative tables in the normalization \eqref{eq:positive-lifts} and
every eligible character component. Successful branches are rechecked
before a certificate is accepted.

None of these functions receives a multiplicity index, historical label,
previous status or positive-list complement as a hypothesis. The
seventeen model matches are a separate positive channel. An arithmetic
failure, field obstruction, missing complete low matrix, exact moment
violation or integer separating vector is a mathematical witness. A
search timeout, unsupported field or unsuccessful linear program is not.
This policy is enforced by the decision ledger.

The top-level allocations in Table~\ref{tab:allocation} are disjoint and
follow the actual decision order. Formal-codegree arithmetic comes first
because it settles almost all inputs without a central search. The
cheap cyclotomic and fusion-subring checks follow. Exact recognition of
cyclic-four and quadratic-parity tensors is cheap; their specialized
systems are built only on those families. Positive reduction and central
induction precede the full signed sphericalization search. These choices
balance rejection frequency and cost, without claiming a formal runtime
optimality theorem. The low-moment bounds are retained for their cost
advantage, despite being consequences of stronger moment conditions.
\begin{table}[ht]\centering\small
\begin{tabular}{lrrrr}\toprule
First accepted top-level certificate&First $\mu$&$\mu\le17$&$\mu\le18$&$\mu\le20$\\\midrule
Formal-codegree $d$-number test&1&23075&28124&40558\\
Cyclotomicity test&5&5&6&8\\
Fusion-subring index test&3&12&12&14\\
Cyclic-four central support&8&3&3&4\\
Quadratic parity and positive induction&3&12&12&20\\
Quadratic parity and spherical characters&12&6&6&6\\
Positive reduction and central induction&2&27&29&30\\
Positive-indicator sphericalization&2&256&282&340\\
\midrule Total general exclusions&&23396&28474&40980\\
\bottomrule\end{tabular}
\caption{Disjoint first-accepted certificates in executable order. First use at a multiplicity does not date the underlying theorem. The order favors cheap general tests, then structural recognition before restricted and full central-induction searches.}\label{tab:allocation}\end{table}

\begin{table}[ht]\centering\small
\setlength{\tabcolsep}{3pt}
\begin{tabular}{r|r|rrrrrrrr|r}\toprule
$\mu$&To decide&D&Cy&Sub&C4&QP&QS&Gate&Sign&Open\\\midrule
1&6&6&0&0&0&0&0&0&0&0\\
2&34&26&0&0&0&0&0&4&4&0\\
3&80&62&0&2&0&2&0&6&8&0\\
4&132&119&0&0&0&0&0&2&11&0\\
5&209&197&2&1&0&0&0&0&9&0\\
6&336&312&0&1&0&0&0&4&19&0\\
7&477&461&0&1&0&0&0&0&15&0\\
8&733&708&0&0&1&0&0&1&23&0\\
9&863&848&1&0&0&0&0&1&13&0\\
10&1082&1064&0&0&0&1&0&2&15&0\\
11&1383&1365&0&1&0&0&0&2&15&0\\
12&1948&1920&0&2&1&1&1&1&22&0\\
13&2211&2188&0&0&0&2&3&0&18&0\\
14&2554&2531&1&1&0&0&1&3&17&0\\
15&2936&2909&1&2&0&1&1&0&22&0\\
16&4157&4125&0&1&1&0&0&1&29&0\\
17&4255&4234&0&0&0&5&0&0&16&0\\
18&5080&5049&1&0&0&0&0&2&26&2\\
19&5640&5609&1&2&0&3&0&0&25&0\\
20&6866&6825&1&0&1&5&0&1&33&0\\
\bottomrule\end{tabular}
\caption{Exact-multiplicity exclusion sequence. To decide omits independently constructed unitary models. Stage symbols follow the row order in Table~\ref{tab:allocation}; each ring appears once, at its first accepted certificate. Open concerns general categorification by these criteria.}\label{tab:stages}\end{table}

\begin{table}[ht]\centering\small
\begin{tabular}{rrrrr}\toprule
$\mu$&Total&Unitary models&General exclusions&General residuals\\\midrule
1&16&10&6&0\\
2&37&3&34&0\\
3&82&2&80&0\\
4&134&2&132&0\\
5&209&0&209&0\\
6&336&0&336&0\\
7&477&0&477&0\\
8&733&0&733&0\\
9&863&0&863&0\\
10&1082&0&1082&0\\
11&1383&0&1383&0\\
12&1948&0&1948&0\\
13&2211&0&2211&0\\
14&2554&0&2554&0\\
15&2936&0&2936&0\\
16&4157&0&4157&0\\
17&4255&0&4255&0\\
18&5080&0&5080&0\\
19&5640&0&5640&0\\
20&6866&0&6866&0\\
\bottomrule\end{tabular}
\caption{Exact automatic coverage of all pinned rank-five tensors through multiplicity twenty. All general, pivotal, spherical, pseudo-unitary and unitary existence questions in this range are decided.}\label{tab:coverage}\end{table}


The last new \emph{criterion family} first used in the checked range is
cyclic-four central support at multiplicity eight: from multiplicity nine
through twenty, every exclusion uses the same repertoire. The
parity-plus-spherical-character \emph{certificate type}, which combines
already available conditions, first occurs at multiplicity twelve.
Thus every top-level certificate type has already appeared by twelve,
and from thirteen through twenty the executable proof allocation
introduces no new type. This identifies the earliest multiplicity in
the verified ledger after which no new criterion family, respectively
certificate type, is needed. It makes no claim of a uniform unbounded
classification: two general residuals occur at eighteen.

For the cyclic-four tensors, automatic recognition returns $k=2,3,4,5$.
The height-six criterion checks respectively $6,6,14,14$ support/indicator
branches. At $k=2$ the balanced Galois test removes two and exact
separators remove four; at $k=3,4,5$ all branches have exact separators.
The same routine leaves the positive $k=1$ model unexcluded.
These are recorded applications of Theorem~\ref{thm:cyclic}, not three
separate nonexistence arguments. Similarly, the twenty-six top-level
quadratic-parity applications are detected from the square rows and
verified modulo two; no fixed list of radicands enters the code.

Integral inputs use the rational specialization of
Theorem~\ref{thm:finite}. No published classification of integral rings
is invoked to exclude them. No one-ring two-moment calculation, special
root-length test or isolated signed-case file is needed. All signed
certificates use the positive-indicator normalization, rather than
retaining separate ordinary and adjoint-parity enumeration modes in the
proof allocation.

\subsection{Proof of Theorem~\ref{thm:main}}
The pinned repository census \cite{Census} gives the counts in
Table~\ref{tab:census}. Section~\ref{sec:models} supplies the
seventeen unitary realizations and their exact tensor matches.
For every other ring with $\mu\le17$, the exhaustive ledger supplies a
general exclusion certificate covered by the theorems and cited inputs
above. There are $23396$ such certificates. At multiplicity eighteen,
$5078$ of the $5080$ rings have general exclusion certificates. At exact
multiplicities nineteen and twenty, every one of the $5640+6866$
repository tensors has a general exclusion certificate. These later
certificates use the same criterion order and exact checker, but are
addressed directly against the pinned repository tensor order, not the
older local basis labels. The two remaining multiplicity-eighteen
rings have complete spherical exclusions. Their bases are
self-dual, so every pivotal structure would be spherical: the left
pivotal dimension of a simple equals the right dimension of its dual.
This excludes pivotal existence as well. The inclusions among the five
existence notions and the seventeen unitary realizations give all the
asserted equivalences and ranges.\qed

The exhaustive step here is a finite, exact certificate verification;
it is not a proof-assistant formalization of the categorical theorems.
The release includes every tensor, branch and integer witness, so its
finite assertions can be reproduced from the pinned census rather than
trusted as a numerical census label. The two proof-ledger verification
commands recheck the entire $28493+12506=40999$-ring source and each
accepted certificate.

\subsection{The two general residuals}
In the twenty self-dual coordinates of Section~\ref{sec:data}, the two
remaining tensors are
\begin{equation}\label{eq:frontier-tensors}\begin{aligned}
v^{(5073)}&=(16,16,17,17,17,17,17,15,17,17,\\ &\hspace{15mm}17,17,17,16,17,18,18,17,17,18),\\
v^{(5076)}&=(17,16,16,17,16,17,17,17,17,17,\\ &\hspace{15mm}17,16,17,17,17,17,17,17,18,18).
\end{aligned}\end{equation}

They have common sorted FP type
\[
 (1,a,a,a+1,a+2),\qquad a=33+2\sqrt{285},
 \qquad D=9120+540\sqrt{285}.
\]
Their formal-codegree polynomials are respectively
\begin{align*}
 (t-3)^2(t-15)(t^2-18240t+68400),\\
 (t-3)(t-5)^2(t^2-18240t+68400).
\end{align*}
Both reduce to $t^2(t+1)^3$ modulo two. Therefore
Lemma~\ref{lem:spectral} restricts a hypothetical sphericalization's
entire dimension field to $\Q(\sqrt{285})$, not merely the field of its
global dimension. However, the displayed FP dimensions do not satisfy
Theorem~\ref{thm:parity}: $285\equiv1\pmod4$ and the irrational
coefficient is two. The applicable ring of integers is
$\ZZ[(1+\sqrt{285})/2]$. Extending the parity rigidity argument to this
situation is a specific possible continuation, not a theorem used here.

At both tensors, a complete spherical certificate was obtained; the
general positive-indicator signed enumeration reached its sixty-second
subprocess limit without completion. The discovery records preserve this
distinction. Thus multiplicity eighteen is the first unresolved
\emph{general} stage of the present pipeline after complete treatment of
sixteen and seventeen. Neither tensor is certified feasible, and neither
is asserted to be unknown throughout the literature.

\subsection{Numbering and reproducibility}
The release's notation $\F{m}{j}$ preserves the first index and basis
from the historical ancillary file of \cite{VS} when one exists.
Classes absent from that file have appended archival labels in local
tensor order. Separately, the field \texttt{local\_index} gives the
one-based ordinal within the archived tensor records at fixed $m$.
At multiplicity eighteen the two indices coincide. Fingerprints
of the full tensors, rather than these labels, key every certificate.

The command \texttt{python verify.py} checks the first $28493$ tensors,
canonical uniqueness, the full pinned-repository isomorphism crosswalk,
and all $28474$ general exclusion
certificates, the seventeen constructed model matches and both spherical
frontier certificates. The companion command
\texttt{python software/verify\_extension.py} checks exact source order,
ring axioms, uniqueness and all $12506$ later general exclusion
certificates. Both regenerate each finite matrix and signed branch.
Python/native Gram comparisons and deliberate certificate mutations are
additional controls. The command \texttt{python -O verify.py} performs
the same proof checks: none is implemented using a removable Python
\texttt{assert}. The accepted separating vectors are tested with integers;
a numerical optimization solver is used only in discovery.

The verification guide gives exact commands, input conventions, checked
integer ranges, dependencies and the failure policy. The bundled pinned
repository tensors are the census input. The discovery command
accepts either an external tensor or a canonical index and
applies the same criteria at the chosen moment height and search budget.
The archive also contains complete decision and provenance ledgers,
input hashes and provenance notes.
